% ------------------------------------------------------------------
%  Capitolo 6 — Le due forze della memoria
%  Parte II
%  Ricerca di riferimento: ricerca 01 §4
%  Voci bibliografiche principali: bjork1992, bjork1994, bjork2011, soderstrom2015,
%    kerr1978, kalyuga2003
%  Epigrafe: sostituita il 25/09/2026 (decisione dell'autore): Aristotele in
%  Diogene Laerzio V, 18 ("radici amare, frutto dolce", O — attribuzione
%  aneddotica, non testo aristotelico diretto; vedi ricerca 03 §2.5 e
%  registro K-12) al posto di Quintiliano XI, 2, 43, spostato
%  come unica epigrafe al cap. 15 (vedi nota di coerenza sotto).
%  Scheda del capitolo: ricerca/06_indice-ragionato.md, capitolo 6
%
%  STATO: prima stesura (25/09/2026), circa 2.700 parole. Note di
%  coerenza qui sotto tutte applicate: lessico identico a quello dei
%  capp. 8-9; "un po' di oblio..." enunciato qui per la prima volta;
%  condizione di superabilità formulata con le parole del cap. 8.
%  Nuova voce bibliografica: kerr1978 (Kerr & Booth, lancio con pratica
%  variata; abstract verificato il 25/09/2026, registro S-53).
%  Nessuna cifra nuova oltre a quelle già verificate.
%
%  NOTE DI COERENZA (rilettura del 25/09/2026, capp. 8-10 -> qui; registro Y-01):
%  - Lessico da riusare identico: "forza di immagazzinamento" e "forza di
%    recupero" (cap. 9, riga ~101): "ogni ricordo ha due forze: quanto e
%    radicato... e quanto e accessibile in questo momento". Il cap. 9 le dà
%    già per note ("quella che abbiamo incontrato nel capitolo sulle due
%    forze"): qui vanno introdotte con questi nomi, non con sinonimi.
%  - Il cap. 9 cita qui il principio "un po' di oblio tra una ripetizione e
%    l'altra rende la ripetizione successiva più efficace" e le "difficoltà
%    desiderabili" di Bjork (bjork2011, già in bibliografia, già citato nel
%    cap. 9). Definirle qui per prime.
%  - Il cap. 8 rimanda qui due volte: la coppia "quanto è radicato / quanto è
%    accessibile" (§ "Perché il recupero funziona") e "la fatica giusta sta
%    nel mezzo: un recupero difficile, ma almeno in parte riuscito" (§
%    "Eliminare la domanda alla prima risposta giusta") — quest'ultima è
%    proprio la condizione di "superabilità" delle difficoltà desiderabili:
%    va enunciata qui in modo da coprire entrambi i rimandi.
%  - RISOLTO (25/09/2026): il doppione di epigrafe (Quintiliano XI, 2, 43 qui
%    e nel cap. 15) è stato sciolto per decisione dell'autore. L'epigrafe di
%    Quintiliano resta solo nel cap. 15 (sonno), dove il passo è più
%    pertinente (parla letteralmente di "una notte interposta"). Qui la
%    nuova epigrafe è Aristotele (radici amare, frutto dolce), proposta già
%    in ricerca 03 §1 e §7 come epigrafe naturale per questo capitolo. Il
%    cap. 9 continua a citare Quintiliano XI, 2, 43 per esteso in nota,
%    rimandando al cap. 15: nessuna modifica necessaria lì.
%  - Nessun rimando diretto dal cap. 10 a questo capitolo (verificato: cap.
%    10 cita solo cap:inganni e cap:tempo, mai cap:dueforze).
%
% ------------------------------------------------------------------
\chapter{Le due forze della memoria}
\label{cap:dueforze}

\epigrafe[Diceva che le radici dell'educazione sono amare, ma il frutto è dolce.]{τῆς παιδείας ἔφη τὰς μὲν ῥίζας εἶναι πικράς, τὸν δὲ καρπὸν γλυκύν.}{Aristotele, in Diogene Laerzio, Vite dei filosofi V, 18}

\iniziale{P}{rova a ricordare} due numeri. Il primo è il numero di
telefono della casa in cui sei cresciuto, o di quella dei tuoi nonni:
quello che da bambino componevi a memoria senza pensarci. Il secondo è
il numero della camera d'albergo in cui hai dormito nell'ultima vacanza,
quello che per una settimana hai ripetuto ogni sera alla reception.

Probabilmente il primo numero non ti viene subito. Ci pensi, ne trovi
le prime cifre, esiti sulle ultime; forse non lo ritrovi affatto. Eppure,
se qualcuno te lo dicesse, lo riconosceresti all'istante, senza ombra di
dubbio, e se dovessi reimpararlo ti basterebbe sentirlo una volta. Il
secondo numero, invece, durante quella settimana ti veniva in mente
senza sforzo: bastava entrare nell'albergo. Oggi, a distanza di qualche
mese, è sparito del tutto, e se te lo dicessero non lo riconosceresti
nemmeno.

Sono due ricordi con proprietà opposte. Il primo è profondamente
radicato ma, in questo momento, difficile da raggiungere. Il secondo era
facilissimo da raggiungere, ma non aveva radici. È un esempio caro a due
psicologi dell'Università della California a Los Angeles, Robert ed
Elizabeth Bjork, che su questa differenza hanno costruito una teoria.
Se questo libro dovesse spiegarne una sola, spiegherebbe questa, perché
tiene insieme quasi tutto ciò che la ricerca ha scoperto su come si
studia: perché il recupero funziona, perché conviene distanziare,
perché mescolare aiuta, e perché tutte queste cose \emph{sembrano}
funzionare peggio di quanto funzionino davvero.

\section{Immagazzinare e ritrovare}

Secondo la teoria dei Bjork, ogni ricordo ha due forze
\parencite{bjork1992}. La prima è quanto è \emph{radicato}: quanto
profondamente è stato appreso, con quanti collegamenti ad altre cose
che sai. La chiamano \emph{forza di immagazzinamento} (\emph{storage
strength}). La seconda è quanto è \emph{accessibile in questo momento}:
quanto facilmente ti viene in mente, adesso, se ti serve. La chiamano
\emph{forza di recupero} (\emph{retrieval strength}). Il numero di casa
dell'infanzia ha un'alta forza di immagazzinamento e una bassa forza di
recupero; il numero della camera d'albergo, durante la vacanza, aveva
un'alta forza di recupero e una bassa forza di immagazzinamento.

Le due forze si comportano in modo diverso. La forza di recupero
oscilla: sale quando usi un ricordo, scende quando passa il tempo o
quando ricordi simili le fanno concorrenza. La forza di
immagazzinamento, secondo la teoria, può soltanto crescere: una volta
che una cosa è stata appresa, non si perde, anche se può diventare
irraggiungibile. È un'affermazione forte, che non si può verificare
direttamente, ma è coerente con molti fatti che conosci già. Ricordi
la distinzione del capitolo precedente tra ciò che è \emph{disponibile}
e ciò che è \emph{accessibile}: le parole che Tulving e Pearlstone
credevano dimenticate tornavano appena si dava l'indizio giusto. E
ricordi il metodo di Ebbinghaus, che misurava il \emph{risparmio} nel
reimparare liste che al momento non sapeva più ripetere
(capitolo~\ref{cap:storia}): se reimparare costa meno che imparare da
capo, qualcosa era rimasto, anche se non si vedeva. Lo sperimenta
chiunque riprenda una lingua studiata anni prima e che credeva di aver
dimenticato: dopo qualche giorno all'estero le parole tornano, più in
fretta di quanto si potesse immaginare.

La teoria dei Bjork si chiama \enquote{nuova teoria del disuso} perché
corregge una teoria più vecchia, secondo la quale i ricordi si
consumano con il disuso, come un sentiero che nessuno percorre più e
che l'erba ricopre. Per i Bjork il disuso non cancella il sentiero:
rende difficile trovarne l'imbocco. Ciò che si perde con il tempo è la
forza di recupero, non la forza di immagazzinamento.

C'è poi un secondo legame tra le due forze, ed è il più utile per chi
studia. Un ricordo con un'alta forza di immagazzinamento non solo
resta, ma si recupera anche più facilmente: quando lo riprendi, la sua
accessibilità risale in fretta e poi cala lentamente. Un ricordo con
poca forza di immagazzinamento, invece, può essere accessibilissimo per
qualche ora, ma la sua accessibilità precipita appena smetti di usarlo.
È esattamente la differenza tra l'esame preparato in tre giorni e
dimenticato in tre settimane, e le cose imparate davvero, che dopo anni
basta sfiorare perché riaffiorino. Lo scopo dello studio, in questi
termini, è semplice da dire: \emph{aumentare la forza di
immagazzinamento}. La forza di recupero del giorno dell'esame verrà di
conseguenza.

\section{Il paradosso: dimenticare aiuta a ricordare}

Ora il punto decisivo. Come si aumenta la forza di immagazzinamento?
Studiando e, soprattutto, recuperando. Ma non ogni volta con lo stesso
guadagno. Secondo i Bjork, \emph{quanto più la forza di recupero è bassa
nel momento in cui studi o richiami qualcosa, tanto più cresce la forza
di immagazzinamento} -- a patto che lo studio o il richiamo riescano.

Pensa a che cosa significa. Se richiami una definizione un minuto dopo
averla letta, la risposta arriva subito: la forza di recupero è ancora
altissima, lo sforzo è nullo, e il guadagno è minimo. Se la richiami il
giorno dopo, quando ha cominciato a sbiadire, devi cercarla; forse ne
ritrovi solo una parte, e il resto lo ricostruisci. È più faticoso, e
proprio per questo il ricordo si radica molto più a fondo. Detto in
modo paradossale: \emph{un po' di oblio tra una ripetizione e l'altra
rende la ripetizione successiva più efficace}.

Da questo paradosso discendono molte delle cose che vedremo nella terza
parte. Ripetere una cosa quando la si sa già perfettamente, sul momento,
è quasi tempo perso; riprenderla quando la si sta dimenticando è il
momento migliore. Per questo conviene richiamare dalla memoria invece
di rileggere, e conviene farlo a distanza di tempo invece che tutto in
una volta. Quintiliano, lo abbiamo visto nel
capitolo~\ref{cap:storia}, aveva osservato che il tempo che di solito
fa dimenticare è lo stesso che rafforza la memoria. Non poteva sapere
perché. La teoria dei Bjork gliene offre una ragione: il tempo abbassa
la forza di recupero, e un recupero riuscito quando la forza di recupero
è bassa fa crescere la forza di immagazzinamento più di qualunque
ripetizione a caldo.\footnote{Quintiliano, \emph{Institutio oratoria}
XI, 2, 43. Il passo, che parla di \enquote{una notte interposta}, sarà
l'epigrafe del capitolo~\ref{cap:corpo}: al paradosso delle due forze
si aggiunge infatti il lavoro di consolidamento che il cervello compie
durante il sonno.}

C'è però una condizione, e non va dimenticata: il recupero deve
\emph{riuscire}, almeno in parte. Se aspetti troppo e non ritrovi più
nulla, il tentativo non ha niente da rafforzare: bisogna ristudiare, e
ricominciare. La fatica utile non è quella di chi cerca a vuoto, ma
quella di chi cerca e, con qualche sforzo, trova.

\section{Le difficoltà desiderabili}

Da qui un'espressione che Robert Bjork ha introdotto negli anni Novanta
e che è diventata una delle più citate della psicologia
dell'apprendimento: \emph{difficoltà desiderabili} (\emph{desirable
difficulties}) \parencite{bjork1994, bjork2011}. Sono le condizioni che
rendono lo studio più lento e faticoso oggi, e il ricordo più saldo e
più flessibile domani. La tabella~\ref{tab:difficolta} ne elenca le
principali, ciascuna accanto al suo contrario \enquote{facile}, quello
che quasi tutti gli studenti preferiscono.

\begin{table}[tbp]
  \centering\small
  \begin{tabular}{@{}>{\raggedright\arraybackslash}p{0.38\textwidth}>{\raggedright\arraybackslash}p{0.32\textwidth}l@{}}
    \toprule
    \textit{Difficoltà desiderabile} & \textit{Il contrario \enquote{facile}} & \textit{Cap.} \\
    \midrule
    Richiamare dalla memoria, a libro chiuso & Rileggere & \ref{cap:recupero} \\
    \addlinespace
    Distanziare le sessioni nel tempo & Concentrare tutto negli ultimi giorni & \ref{cap:tempo} \\
    \addlinespace
    Alternare tipi diversi di problemi & Esercitarsi su un tipo alla volta & \ref{cap:alternanza} \\
    \addlinespace
    Generare la risposta prima di vederla & Ricevere la risposta pronta & \ref{cap:capire} \\
    \addlinespace
    Variare esempi, condizioni e contesti & Studiare sempre nello stesso modo & \ref{cap:tempo}, \ref{cap:alternanza} \\
    \bottomrule
  \end{tabular}
  \caption{Le principali difficoltà desiderabili e il loro contrario
    \enquote{facile}. In ogni riga la colonna di sinistra rallenta lo
    studio e migliora il ricordo a lungo termine; quella di destra fa il
    contrario.}
  \label{tab:difficolta}
\end{table}

Tutte le voci della tabella hanno qualcosa in comune: costringono la
mente a lavorare quando lavorare costa. Richiamare costringe a cercare
ciò che rileggendo si trova già pronto. Distanziare fa sì che, a ogni
ripresa, il ricordo sia un po' sbiadito e vada ricostruito. Alternare
obbliga a chiedersi ogni volta di che problema si tratti, invece di
ripetere la procedura dell'esercizio precedente. Generare, cioè provare
a rispondere prima di conoscere la risposta, attiva ciò che sai e
prepara un posto per ciò che imparerai.

L'ultima riga merita un esempio, perché è la meno intuitiva. Alla fine
degli anni Settanta due ricercatori, Robert Kerr e Bernard
Booth, allenarono alcuni bambini di otto anni a lanciare un sacchetto
di sabbia su un bersaglio posto a terra, senza vederlo al momento del
lancio \parencite{kerr1978}. Un gruppo si esercitava sempre alla
distanza su cui sarebbe stato poi messo alla prova; l'altro non si
esercitava mai a quella distanza, ma a due distanze diverse, una più
corta e una più lunga. Alla prova finale, sulla distanza che i primi
avevano provato e riprovato, andarono meglio i secondi. Variare la
pratica aveva reso l'abilità più flessibile, più capace di adattarsi a
una situazione nuova, anche se durante l'allenamento sembrava meno
mirata. Lo stesso vale per lo studio: chi incontra un concetto in
esempi diversi, formulato in modi diversi, applicato a casi diversi,
lo riconosce e lo usa meglio di chi lo ha sempre visto nella stessa
forma.

\subsection*{Desiderabili, cioè superabili}

Una difficoltà, però, è desiderabile \emph{solo se riesci a superarla}.
È la condizione che abbiamo già incontrato a proposito del recupero, ed
è quella che più spesso si dimentica. La fatica giusta sta nel mezzo:
un recupero difficile, ma almeno in parte riuscito; un problema che ti
costringe a pensare, ma che con un po' di sforzo, e magari con un aiuto,
sai risolvere. Se la difficoltà supera le tue forze, perché ti mancano
le conoscenze di base, perché il materiale è troppo nuovo o troppo
complesso, non costruisce nulla: produce soltanto frustrazione, e un
sovraccarico della memoria di lavoro che conosci dal
capitolo~\ref{cap:cervello}.

Ne derivano due avvertenze. La prima: non ogni difficoltà è
desiderabile. Un testo stampato male, un'aula rumorosa, una lezione
troppo veloce rendono lo studio più faticoso, ma non attivano nessuno
dei processi che radicano i ricordi: sono difficoltà e basta. La
seconda: la dose giusta dipende da quanto sai. Chi comincia ha bisogno
di più guida -- spiegazioni chiare, esempi svolti, domande più facili
-- e di meno difficoltà; man mano che le conoscenze crescono, la guida
va ridotta e le difficoltà aumentate, perché ciò che aiuta il
principiante può diventare inutile, o perfino dannoso, per chi è già
esperto \parencite{kalyuga2003}. Una difficoltà desiderabile non è una
prova di coraggio: è una fatica misurata sulle tue forze.

\begin{inpratica}
\textbf{Come riconoscere la fatica giusta.} Mentre studi, chiediti in
quale di queste tre situazioni ti trovi.
\begin{itemize}
  \item \emph{Troppo facile.} Rispondi all'istante, senza pensarci,
        perché hai appena letto la risposta. Non è tempo perso del
        tutto, ma rende poco: aspetta di più prima di interrogarti,
        mescola le domande, cambia la formulazione.
  \item \emph{Giusto.} Devi cercare; ritrovi una parte, a volte
        sbagli, poi controlli e correggi. È qui che la forza di
        immagazzinamento cresce di più. Resta in questa zona, anche se
        è scomoda.
  \item \emph{Troppo difficile.} Non ricordi quasi nulla, o non capisci
        nemmeno la domanda. Non insistere a vuoto: torna al testo, a un
        esempio svolto, a una spiegazione più semplice; poi riprova,
        magari il giorno dopo.
\end{itemize}
Un buon segnale pratico: se in una sessione di recupero non sbagli mai,
stai lavorando troppo sul facile; se sbagli quasi sempre, ti manca
ancora la base. La zona utile è quella in cui riesci, con fatica, più
spesso di quanto fallisci.
\end{inpratica}

\section{Prestazione non è apprendimento}

C'è un'ultima conseguenza della teoria, e forse è la più importante per
chi studia e per chi insegna. Se lo scopo è aumentare la forza di
immagazzinamento, e se ciò che vediamo mentre studiamo è soprattutto la
forza di recupero, allora \emph{ciò che vediamo mentre studiamo non ci
dice quanto stiamo imparando}.

Nicholas Soderstrom e Robert Bjork hanno raccolto decenni di ricerche
su questo punto, e lo hanno formulato come una distinzione tra due
cose che il linguaggio comune confonde \parencite{soderstrom2015}. La
\emph{prestazione} è ciò che si osserva durante lo studio o subito
dopo: quanti esercizi riesci a fare, quanto velocemente rispondi, quanto
ti sembra chiaro ciò che hai letto. L'\emph{apprendimento} è il
cambiamento duraturo che resta quando la prestazione del momento è
svanita: ciò che saprai fare tra un mese, all'esame, nella vita. Le due
cose possono andare in direzioni opposte.

Può esserci prestazione senza apprendimento. Rileggere tre volte lo
stesso capitolo la sera prima dell'esame produce una prestazione
brillante il mattino dopo, e un apprendimento modesto: la forza di
recupero è alta, quella di immagazzinamento è cresciuta poco. Lo
stesso vale per gli esercizi fatti tutti dello stesso tipo, uno dopo
l'altro: il decimo viene meglio del primo, ma solo perché la procedura
è ancora lì, calda, dall'esercizio precedente. E può esserci
apprendimento senza prestazione: le difficoltà desiderabili rallentano
lo studio, moltiplicano gli errori, danno la sensazione di non
progredire, mentre la forza di immagazzinamento cresce in silenzio.
Negli esperimenti di questo libro lo vedrai accadere più volte: cinque
minuti dopo lo studio vince chi ha riletto, una settimana dopo vince
chi si è interrogato (capitolo~\ref{cap:recupero}); durante la pratica
vanno meglio i blocchi, alla prova finale vince l'alternanza
(capitolo~\ref{cap:alternanza}).

Per chi studia, la conclusione è scomoda ma liberatoria. Non puoi
giudicare il tuo metodo dalla sensazione di scorrevolezza che provi
mentre lo usi, né dai risultati che ottieni nel corso della stessa
sessione. Un'ora di studio piena di vuoti e di errori può averti
insegnato molto più di un'ora in cui tutto filava liscio. Per chi
insegna, la conclusione è la stessa vista dall'altra parte della
cattedra: una lezione in cui tutti annuiscono, un'esercitazione in cui
tutti rispondono bene al primo colpo, non dimostrano che gli studenti
hanno imparato; possono dimostrare soltanto che le risposte erano
ancora fresche.

Il detto che la tradizione antica attribuiva ad Aristotele, e che apre
questo capitolo, lo diceva con un'immagine: le radici dell'educazione
sono amare, ma il frutto è dolce.\footnote{Diogene Laerzio, \emph{Vite
dei filosofi} V, 18. È un detto attribuito ad Aristotele dalla
tradizione aneddotica, non una frase delle sue opere.} Non tutte le
amarezze, però, portano frutto: soltanto quelle che fanno lavorare la
mente nel modo giusto, e che si riescono a superare. Resta un problema.
Se le strategie migliori sono anche quelle che sembrano peggiori, come
può uno studente accorgersene? Nella maggior parte dei casi non se ne
accorge, e sceglie le altre. Perché accada, e come evitarlo, è
l'argomento del prossimo capitolo.

\begin{daricordare}
La fatica giusta è un buon segno, purché tu riesca, almeno in parte. Ciò
che è facile mentre studi misura la forza di recupero del momento, non
ciò che resterà.
\end{daricordare}

\begin{domande}
  \item Qual è la differenza tra forza di immagazzinamento e forza di
        recupero? Fai un esempio di ricordo con molta dell'una e poca
        dell'altra.
  \item Perché un recupero difficile rafforza la memoria più di uno
        facile? Che cosa succede se il recupero non riesce affatto?
  \item Perché una difficoltà deve essere \enquote{superabile} per
        essere desiderabile? Fai un esempio di difficoltà che non lo è.
  \item Scegli un esame che stai preparando: quali difficoltà
        desiderabili potresti introdurre nel tuo modo di studiarlo?
  \item Qual è la differenza tra prestazione e apprendimento? Perché la
        prestazione durante lo studio può ingannare?
  \item (Ripresa, capitolo~\ref{cap:cervello}) Quanto è grande la
        memoria di lavoro? Che cosa c'entra con le difficoltà che non
        sono desiderabili?
  \item (Ripresa, capitolo~\ref{cap:storia}) Che cosa misurava
        Ebbinghaus con il metodo del \enquote{risparmio}, e che cosa
        dice sulla forza di immagazzinamento?
\end{domande}

\begin{sintesi}
  \item Ogni ricordo ha due forze: quanto è radicato (forza di
        immagazzinamento) e quanto è accessibile in questo momento
        (forza di recupero). Il tempo fa calare la seconda, non la
        prima.
  \item Il paradosso: quanto più un ricordo è difficile da raggiungere,
        tanto più un recupero riuscito lo rafforza. Un po' di oblio tra
        una ripetizione e l'altra rende più efficace la ripetizione
        successiva.
  \item Le difficoltà desiderabili (richiamare, distanziare, alternare,
        generare, variare) rallentano lo studio e migliorano il
        ricordo, a condizione di essere superabili.
  \item Prestazione e apprendimento sono cose diverse: ciò che si vede
        mentre si studia non è ciò che resta.
\end{sintesi}

\finecapitolo
